\documentclass[10pt,a4paper]{amsart}
\usepackage{amsmath,amssymb,amsthm,graphicx}
% Retained minimal preamble for the Stacks Project chapter more-algebra.tex (More on Algebra).
% The chapter file pinned-original-source.tex carries only \input{preamble} on its first line, so no
% preamble text lives inside the chapter itself.  The Stacks Project uses one single shared preamble for
% every chapter; the snapshot stages it identically as sources/stacks-chow/preamble.tex and
% sources/stacks-etale-cohomology/preamble.tex.  The block below is copied mechanically, line for line,
% from sources/stacks-chow/preamble.tex lines 14-18 (the Xy-pic setup) and 154-217 (the native theorem
% declarations together with the mathematical macros).  Nothing was added, removed or redefined; the
% wrapper class and the amsmath/amssymb/amsthm/graphicx base packages are supplied by the assembler.
% Chapter: more-algebra.tex -- More on Algebra.
% For commutative diagrams we use Xy-pic
\usepackage[all]{xy}

% We use 2cell for 2-commutative diagrams.
\xyoption{2cell}

% Theorem environments.
%
\theoremstyle{plain}
\newtheorem{theorem}[subsection]{Theorem}
\newtheorem{proposition}[subsection]{Proposition}
\newtheorem{lemma}[subsection]{Lemma}

\theoremstyle{definition}
\newtheorem{definition}[subsection]{Definition}
\newtheorem{example}[subsection]{Example}
\newtheorem{exercise}[subsection]{Exercise}
\newtheorem{situation}[subsection]{Situation}

\theoremstyle{remark}
\newtheorem{remark}[subsection]{Remark}
\newtheorem{remarks}[subsection]{Remarks}

\numberwithin{equation}{subsection}

% Macros
%
\def\lim{\mathop{\mathrm{lim}}\nolimits}
\def\colim{\mathop{\mathrm{colim}}\nolimits}
\def\Spec{\mathop{\mathrm{Spec}}}
\def\Hom{\mathop{\mathrm{Hom}}\nolimits}
\def\Ext{\mathop{\mathrm{Ext}}\nolimits}
\def\SheafHom{\mathop{\mathcal{H}\!\mathit{om}}\nolimits}
\def\SheafExt{\mathop{\mathcal{E}\!\mathit{xt}}\nolimits}
\def\Sch{\mathit{Sch}}
\def\Mor{\mathop{\mathrm{Mor}}\nolimits}
\def\Ob{\mathop{\mathrm{Ob}}\nolimits}
\def\Sh{\mathop{\mathit{Sh}}\nolimits}
\def\NL{\mathop{N\!L}\nolimits}
\def\CH{\mathop{\mathrm{CH}}\nolimits}
\def\proetale{{pro\text{-}\acute{e}tale}}
\def\etale{{\acute{e}tale}}
\def\QCoh{\mathit{QCoh}}
\def\Ker{\mathop{\mathrm{Ker}}}
\def\Im{\mathop{\mathrm{Im}}}
\def\Coker{\mathop{\mathrm{Coker}}}
\def\Coim{\mathop{\mathrm{Coim}}}

% Boxtimes
%
\DeclareMathSymbol{\boxtimes}{\mathbin}{AMSa}{"02}

%
% Macros for moduli stacks/spaces
%
\def\QCohstack{\mathcal{QC}\!\mathit{oh}}
\def\Cohstack{\mathcal{C}\!\mathit{oh}}
\def\Spacesstack{\mathcal{S}\!\mathit{paces}}
\def\Quotfunctor{\mathrm{Quot}}
\def\Hilbfunctor{\mathrm{Hilb}}
\def\Curvesstack{\mathcal{C}\!\mathit{urves}}
\def\Polarizedstack{\mathcal{P}\!\mathit{olarized}}
\def\Complexesstack{\mathcal{C}\!\mathit{omplexes}}
% \Pic is the operator that assigns to X its picard group, usage \Pic(X)
% \Picardstack_{X/B} denotes the Picard stack of X over B
% \Picardfunctor_{X/B} denotes the Picard functor of X over B
\def\Pic{\mathop{\mathrm{Pic}}\nolimits}
\def\Picardstack{\mathcal{P}\!\mathit{ic}}
\def\Picardfunctor{\mathrm{Pic}}
\def\Deformationcategory{\mathcal{D}\!\mathit{ef}}

% --- appended 2026-10-12 (Chief): the project own definitions of the reference/slogan/history
% display environments, copied mechanically from the shared Stacks preamble lines 8-12
% (\usepackage{verbatim} plus the three \newenvironment lines, which the chapter preamble keeps
% verbatim from the project preamble). Nothing added, removed or redefined: these environments
% carry no mathematical content and the project itself typesets them as comments. Without these
% three lines a counted statement whose own text uses them fails with Environment slogan undefined.
\usepackage{verbatim}
\newenvironment{reference}{\comment}{\endcomment}
\newenvironment{slogan}{\comment}{\endcomment}
\newenvironment{history}{\comment}{\endcomment}
\usepackage{cleveref}
\crefname{originallemma}{original lemma}{original lemmas}
\Crefname{originallemma}{Original lemma}{Original lemmas}
\makeatletter
\newlabel{algebra-lemma-hilbert-ses-chi}{{00KC}{source}{}{sourcepointer-1}{}}
\newlabel{algebra-lemma-hilbert-ses-chi@cref}{{[originallemma][0][]00KC}{[1][0][]source}{}{}{}}
\makeatother
\crefname{originallemma}{original lemma}{original lemmas}
\Crefname{originallemma}{Original lemma}{Original lemmas}
\makeatletter
\newlabel{algebra-lemma-support-dimension-d}{{00L8}{source}{}{sourcepointer-2}{}}
\newlabel{algebra-lemma-support-dimension-d@cref}{{[originallemma][0][]00L8}{[1][0][]source}{}{}{}}
\makeatother
\crefname{originallemma}{original lemma}{original lemmas}
\Crefname{originallemma}{Original lemma}{Original lemmas}
\makeatletter
\newlabel{lemma-minprimespoly}{{38016}{source}{}{sourcepointer-3}{}}
\newlabel{lemma-minprimespoly@cref}{{[originallemma][0][]38016}{[1][0][]source}{}{}{}}
\makeatother
\crefname{originallemma}{original lemma}{original lemmas}
\Crefname{originallemma}{Original lemma}{Original lemmas}
\makeatletter
\newlabel{lemma-multiplicity}{{38261}{source}{}{sourcepointer-4}{}}
\newlabel{lemma-multiplicity@cref}{{[originallemma][0][]38261}{[1][0][]source}{}{}{}}
\makeatother
\title{The number of top-dimensional minimal primes of a Noetherian local ring is bounded by the length of the quotient by a system of parameters}
\author{The Stacks Project Authors}
\date{stacks-project:more-algebra.tex:local-snapshot-20261012; GNU FDL 1.2 or later}
\begin{document}
\maketitle
\setcounter{subsection}{1}
\noindent\textit{Editorial note.} All delivered mathematical text is the author's own, copied from the pinned Stacks Project chapter file more-algebra.tex; no statement or proof step was rewritten, reordered or removed.\par\medskip
\noindent\textit{Editorial note.} The chapter file pinned-original-source.tex carries only \textbackslash{}input\{preamble\} on its first line, so no preamble text lives inside the chapter. The retained file sources/stacks-more-algebra/retained-preamble.tex is a mechanical, line-for-line copy of the project's single shared preamble as staged in the snapshot; nothing was added, removed or redefined and no package was added.\par\medskip
\noindent\textit{Editorial note.} The source URL is the source-level Stacks Project URL, as required when more than one card is produced from one source; the per-claim tag is carried by claim\_key.\par\medskip
\noindent\textit{Editorial note.} The statement writes R/(g\_1, ..., g\_n) although the system of parameters has d entries, so the symbol n does not denote the dimension; that is the author's text and is reproduced unchanged. The author also writes 'of degree exactly(!) d' inside the counted proof span, reproduced unchanged. The six references inside the counted proof span are recorded as non-load-bearing navigation pointers.\par\medskip
\noindent\textit{Allowed standard prerequisites (editorial).}\begin{itemize}
\item Hilbert functions of an ideal in a local ring and the degree of the polynomial giving the length of a quotient. used at proof lines 38330-38343. Is the asymptotic tool that turns the support computation into a numerical bound.
\item Length additivity in a short exact sequence of finite modules over a local ring. invoked at proof lines 38333-38338. Is the comparison that relates the three Hilbert polynomials.
\item The multiplicity lemma bound relating the leading coefficient to the length. invoked at proof lines 38344. Converts the leading-coefficient lower bound t/d! into the asserted length inequality.
\item The one-dimensional case of the statement. invoked at proof line 38306. Is the base case of the reduction.
\end{itemize}\medskip
\ifdefined\hypertarget\hypertarget{sourcepointer-1}{}\fi
\noindent\textit{Source reference.} Original lemma 00KC: focused-seeds/AGP/sources/stacks-snapshot-20261011/tags.txt, source lines 746--746. Cross-chapter navigation pointer into the Stacks Project; the standard fact it supplies is applied inside the counted proof span and the argument resting on it stays there.\par\medskip
\ifdefined\hypertarget\hypertarget{sourcepointer-2}{}\fi
\noindent\textit{Source reference.} Original lemma 00L8: focused-seeds/AGP/sources/stacks-snapshot-20261011/tags.txt, source lines 778--778. Cross-chapter navigation pointer into the Stacks Project; the standard fact it supplies is applied inside the counted proof span and the argument resting on it stays there.\par\medskip
\ifdefined\hypertarget\hypertarget{sourcepointer-3}{}\fi
\noindent\textit{Source reference.} Original lemma 38016: sources/stacks-more-algebra/pinned-original-source.tex, source lines 38016--38016. Same-chapter navigation pointer; the statement it names is standard for the section and the argument that consumes it stays inside the counted proof span.\par\medskip
\ifdefined\hypertarget\hypertarget{sourcepointer-4}{}\fi
\noindent\textit{Source reference.} Original lemma 38261: sources/stacks-more-algebra/pinned-original-source.tex, source lines 38261--38261. Same-chapter navigation pointer; the statement it names is standard for the section and the argument that consumes it stays inside the counted proof span.\par\medskip
\begin{lemma}
\label{lemma-minprimespolyhigher}
Let $(R, \mathfrak m)$ be a Noetherian local ring of dimension $d$, let $t$
be the number of minimal prime ideals of $R$ of dimension $d$, and let
$(g_1,\ldots,g_d)$ be a system of parameters. Then
$t \leq \text{length}_R(R/(g_1,\ldots,g_n))$.
\end{lemma}

\begin{proof}
If $d = 0$ the lemma is trivial. If $d = 1$ the lemma is
Lemma \ref{lemma-minprimespoly}. Thus we may assume $d > 1$.
Let $\mathfrak p_1, \ldots, \mathfrak p_s$ be the minimal prime ideals of
$R$ where the first $t$ have dimension $d$, and denote
$I = (g_1, \ldots, g_n)$. Arguing in exactly the same way as in
the proof of Lemma \ref{lemma-minprimespoly} we can assume $R$ is reduced.

\medskip\noindent
Assume $R$ is reduced with minimal primes
$\mathfrak p_1, \ldots, \mathfrak p_t$.
This means there is an exact sequence
$$
0 \to R \to
\prod\nolimits_{i = 1}^t R/\mathfrak p_i \to Q \to 0
$$
Here $Q$ is the cokernel of the first map.
Write $M = \prod_{i = 1}^t R/\mathfrak p_i$.
Localizing at $\mathfrak p_j$ we see that
$$
R_{\mathfrak p_j} \to M_{\mathfrak p_j} =
\left(\prod\nolimits_{i=1}^t R/\mathfrak p_i\right)_{\mathfrak p_j} =
(R/\mathfrak p_j)_{\mathfrak p_j}
$$
is surjective. Thus $Q_{\mathfrak p_j} = 0$ for all $j$. Therefore no
height $0$ prime of $R$ is in the support of $Q$. It follows that
the degree of the numerical polynomial
$n \mapsto \text{length}_R(Q/I^nQ)$ equals $\dim(\text{Supp}(Q)) < d$, see
Algebra, Lemma \ref{algebra-lemma-support-dimension-d}.
By Algebra, Lemma \ref{algebra-lemma-hilbert-ses-chi}
(which applies as $R$ does not have finite length) the polynomial
$$
n \longmapsto
\text{length}_R(M/I^nM) - \text{length}_R(R/I^n) - \text{length}_R(Q/I^nQ)
$$
has degree $< d$. Since $M = \prod R/\mathfrak p_i$ and since
$n \to \text{length}_R(R/\mathfrak p_i + I^n)$ is a numerical
polynomial of degree exactly(!) $d$ for $i = 1, \ldots, t$ (by
Algebra, Lemma \ref{algebra-lemma-support-dimension-d})
we see that the leading coefficient of $n \mapsto \text{length}_R(M/I^nM)$
is at least $t/d!$. Thus we conclude by Lemma \ref{lemma-multiplicity}.
\end{proof}

\end{document}
